\documentclass[10pt,a4paper]{amsart}
\usepackage[margin=19mm]{geometry}
\usepackage{amsmath,amssymb,mathtools}
\usepackage[scr=rsfs,cal=euler]{mathalfa}
\usepackage{xfrac}
\usepackage[unicode,hidelinks,bookmarks=false]{hyperref}
\newcommand{\ie}{i.e.\ }
\newcommand{\cf}{cf.\ }
\newcommand{\resp}{respectively\ }
\newcommand{\Subsection}{\subsection}
\providecommand{\nobreakdash}{\nobreak\hbox{-}}
\setlength{\emergencystretch}{2em}
\multlinegap0pt
\usepackage{bm}
\usepackage{bbm}
\usepackage{dsfont}
\usepackage{xspace}
\usepackage{cancel}
\usepackage{mathtools}
\usepackage{amsthm}
\makeatletter
\@ifundefined{definition}{\newtheorem{definition}{Definition}}{}
\@ifundefined{lemma}{\newtheorem{lemma}[definition]{Lemma}}{}
\makeatother
\DeclareMathOperator{\id}{id}
\title[Coarse Homotopy Theory and Shape Theory]{Coarse Homotopy Theory and Shape Theory}
\author{Felix Lange}
\date{Source: arXiv:2606.14212v1 (CC BY 4.0)}
\begin{document}
\maketitle

\noindent\emph{Source and licence.} Coarse Homotopy Theory and Shape Theory. Version arXiv:2606.14212v1 (CC BY 4.0), distributed under the Creative Commons Attribution 4.0 International licence, as recorded in the arXiv licence metadata for this exact version and shown on the abstract page https://arxiv.org/abs/2606.14212v1. Full rights notice: licenses/arxiv-2606.14212-CC-BY-4.0.txt.\par\smallskip

\noindent\textit{Editorial scope.} Counted unit: the Lemma whose source label is \texttt{che} in the article (source0.tex lines 5675--5676, source label \texttt{che}), delivered together with the article's own complete original proof of it (source0.tex lines 5678--5832, one \texttt{proof} environment of 155 lines). No mathematics is written, repaired or re-derived: statement and proof are the article's own text line for line. Required context retained uncounted: shrinkindep (lines 2020--2021); shrinkcoarse (lines 5480--5481); globallycoarse (lines 5627--5633). Every definition, display and label of those lines is retained unchanged. Omitted from this package and not redistributed: 1--2019, 2022--5479, 5482--5626, 5634--5674, 5677--5677, 5833--7510. Nothing inside a retained excerpt is omitted or altered: every retained body is delivered exactly as the article's lines are written, so no omission receipt of any kind accompanies this package. Citations: the retained excerpts carry no citation command at all, so no bibliography entry of the article is transcribed and no reference is redistributed. Reference adaptation: none was needed and none was made; every reference inside the delivered unit resolves inside it, so no unresolved reference or citation is produced. Printed slips kept literal: no printed word, symbol, hypothesis or number of the counted statement or of its proof was altered. Layout and class: the article's own class is \texttt{article} with packages including geometry, inputenc, fontenc, babel, lmodern, csquotes, enumitem, float, hyperref, amsthm, amsmath, amssymb, amsfonts, mathrsfs; the wrapper is the lane's pinned \texttt{amsart} editorial wrapper at 10pt on A4, which restates the theorem-like environments the retained text needs and adds the article's own macro definitions that the retained text needs (\texttt{\textbackslash id}), with extra wrapper packages (bm, bbm, dsfont, xspace, cancel, mathtools). The delivered numbering is the unit's own. Assets: no retained span contains a figure environment, a graphics-inclusion command, a PDF-page inclusion, a table or a TikZ picture; no image file is copied into this package, the package declares no asset dependency and no image file is redistributed with it. Licence: version arXiv:2606.14212v1 of this article is distributed under the Creative Commons Attribution 4.0 International licence, as recorded in the arXiv licence metadata for this exact version and shown on the abstract page https://arxiv.org/abs/2606.14212v1. A full rights notice is delivered at licenses/arxiv-2606.14212-CC-BY-4.0.txt. Characters: every retained excerpt is pure ASCII, so the delivered engine, XeLaTeX with its default Latin Modern text font, reports no missing character.
\begin{lemma} \label{shrinkindep} Let $s_{\gamma},s_{\gamma'}$ be two shrinking maps. There is a linear coarse homotopy $G: I_{p}(cX)\rightarrow cX$ between $s_{\gamma}$ and $s_{\gamma'}$.
\end{lemma}

\begin{lemma} The map $i$ can be made coarse, by precomposing with a shrinking map $s$ on $cX$. \label{shrinkcoarse}
\end{lemma}

\begin{lemma} \begin{enumerate} \label{globallycoarse}
\item There exists a coarse shrinking map $S_{\xi'}$ on $M_{Y}$, defined as the composition $S_{\xi'}:=i' \circ s_{\xi'} \circ R$, such that $\xi' S_{\xi'}$ is coarse. 
\item There exists a coarse shrinking map $S_{\psi \xi'}$ on $I_{p}(M_{Y})$, defined as the composition $S_{\psi \xi'}:= (i' \times \id) \circ s_{\psi \xi'}  \circ (R \times \id)$, where $s_{\psi \xi'}: I_{p}(cY)\rightarrow I_{p}(cY)$ is a shrinking map on $I_{p}(cY)$, such that $F S_{\psi \xi'}$ is coarse. 
\item There exists a coarse shrinking map $S_{\xi'\psi}$ on $I_{p}(M_{X})$, defined as the composition $S_{\xi'\psi}:=(i' \times \id) \circ s_{\xi'\psi}  \circ (R \times \id)$, where $s_{\xi' \psi}: I_{p}(cX)\rightarrow I_{p}(cX)$ is a shrinking map on $I_{p}(cY)$, such that $G S_{\xi' \psi}$ is coarse. 
\item $S_{\xi'}, S_{\psi \xi'}$ and $S_{\xi'\psi}$ are coarsely homotopic to the identity. 
\end{enumerate}
\end{lemma}

\begin{lemma} \label{che} $\xi' S_{\xi'}$ is a coarse homotopy inverse to $is_{l}\mathfrak{a}R: M_{X}\rightarrow M_{Y}$.  
\end{lemma}

\begin{proof} We compute both compositions. Recall that $\psi = \widetilde{is_{l}\mathfrak{a}R}$ for sufficiently large $l$. By Lemma \ref{globallycoarse}, $F S_{\psi \xi'}$ gives us a coarse homotopy between $\psi \xi' S_{\psi \xi'}$ and $\id_{M_{Y}}S_{\psi \xi'}$. (Note that since the shrinking map $S_{\psi \xi'}$ is uniform in time we have used the same symbol for its restriction to $t=0,1$.) We wish to show the homotopy $\psi \xi' S_{\xi'} \simeq \psi \xi' S_{\psi \xi'}$. This is a similar calculation to Lemma \ref{shrinkindep}. Let $\gamma: [1,\infty)\rightarrow [1,\infty)$ be the shrinking map associated to $s_{\xi'}$ and let $\gamma'$ be associated to $s_{\psi\xi'}$. We can assume that $\gamma'(t)\leq \gamma(t)$ for all $t\in [1,\infty)$. Let 
\begin{align*}
H: I_{p}(cY)&\rightarrow cY\\
(nx,n,nt)&\mapsto (\eta(n,t)x,\eta(n,t))
\end{align*}
be the coarse homotopy from $s_{\xi'}$ to $s_{\psi\xi'}$, where $\eta(n,t) = (1-t)\gamma(n)+t\gamma'(n)$.  We show that $\xi' i' H$ is uniformly locally Lipschitz. Let $\{m_{h}\}_{h\in \mathbb{N}_0}$ be the sequence of constants associated to $\gamma$. Consider two points $(nx,n,nt),(sy,s,sl)$ in $c_{2^{m_{h}}(2^{h},2^{h+1}]}(Y \times [0,1])$. 
\begin{align*}
d_{cY}(H(nx,n,&nt),H(sy,s,sl)) \leq \|\eta(n,t)x - \eta(n,t)y\| + 2|\eta(n,t)-\eta(s,l)|\\
&\leq \frac{n}{2^{m_{h}}} \|x-y\| + 2(|\eta(n,t)-\eta(s,t)|+|\eta(s,t)-\eta(s,l)|)\\
&\leq \frac{1}{2^{m_{h}}} \|nx-ny\|+2(\frac{1}{2^{m_{h}}}|n-s|+|t-l||\gamma(s)+\gamma'(s)|)\\
&\leq \frac{1}{2^{m_{h}}} (\|nx-sy\| + |n-s|) + 2(\frac{1}{2^{m_{h}}}|n-s| +\frac{2}{2^{m_{h}}}|st-sl|)\\
&\leq \frac{1}{2^{m_{h}}} (\|nx-sy\| + |n-s|) + 2(\frac{1}{2^{m_{h}}}|n-s| + \frac{2}{2^{m_{h}}}(|nt-sl|+|n-s|))\\
&= \frac{1}{2^{m_{h}}}(\|nx-sy\| + 7|n-s| + 4|nt-sl|)
\end{align*}
The images of the points have heights in $[2^h,2^{h+1}]$. Therefore the restriction of $\xi' i' H$ to $c_{2^{m_{h}}(2^{h},2^{h+1}]}(Y \times [0,1])$ is $7$-Lipschitz. For points in the height interval $(2^{h+m_{h-1}},2^{h+m_{h}}]$ the calculation is analagous. We have  
\begin{align*}
d_{cY}(H(nx,n,nt),H(sy,s,sl)) \leq \frac{1}{2^{m_{h-1}}}(\|nx-sy\| + 7|n-s| + 4|nt-sl|)
\end{align*}
with image at height $2^{h}$. Therefore $\xi' i' H$ restricted to $c_{(2^{h+m_{h-1}},2^{h+m_{h}}]}(Y\times [0,1])$ is also $7$-Lipschitz. This shows that $\xi' i' H$ is uniformly locally $7$-Lipschitz, and therefore coarse. Hence the composition $\psi \circ (\xi' i' H)\circ (R\times \id)$ is a coarse homotopy between $\psi \xi' S_{\xi'}$ and $\psi \xi' S_{\psi \xi'}$. We have the following sequence of coarse homotopies:
\begin{align*}
(is_{l}\mathfrak{a}R) \xi' S_{\xi'} \simeq \psi \xi' S_{\xi'} \simeq \psi \xi' S_{\psi \xi'} \simeq \id_{M_{Y}} S_{\psi \xi'} \simeq \id_{M_{Y}}
\end{align*}
where in the first equivalence we have used that $\psi$ is close to $is_{l}\mathfrak{a} R$.\\

 To show that $ \xi' S_{\xi'} (is_{l}\mathfrak{a}R) \simeq \id_{M_{X}}$ is more complicated. 
We have a coarse homotopy $GS_{\xi' \psi}$ between $\xi' \psi S_{\xi' \psi}$ and $\id_{M_{X}} S_{\xi' \psi}$. We would like to have the following sequence of homotopies:
\begin{align*}
 \xi' S_{\xi'} (is_{l}\mathfrak{a}R) \simeq   \xi' S_{\xi'} \psi \simeq \xi' S_{\xi'} \psi S_{\xi' \psi} \simeq \xi' \psi S_{\xi' \psi} \simeq \id_{M_{X}} S_{\xi' \psi} \simeq \id_{M_{X}}
\end{align*}
The trouble appears in the homotopy $ \xi' S_{\xi'} \psi S_{\xi' \psi} \simeq \xi' \psi S_{\xi' \psi}$. Since $\xi'$ is only locally Lipschitz, we cannot simply run the homotopy $\xi' H \psi S_{\xi' \psi}$ where $H$ is a coarse homotopy between $S_{\xi'}$ and $\id$ without thinking: any bounded distance can become arbitrarily large under $\xi'$. We therefore run the homotopy in two steps:\\

Let $\gamma$ be the shrinking map corresponding to $s_{\xi'}$, $\gamma'$ the shrinking map corresponding to $s_{\xi' \psi}$. We let $\mathbf{D}$ be the deformation retract
\begin{align*}
\mathbf{D}: I_{p}(M_{Y})&\rightarrow M_{Y}\\
(x,n,nt)&\mapsto \mathfrak{q}_{(1-t)n + t\gamma(n)}(x,n)
\end{align*}
where $\mathfrak{q}$ refers to the height projection. $\mathbf{D}_{1}: M_{Y} \rightarrow M_{Y}$ is the map which shrinks heights via $\gamma$. We first compute the Lipschitz constant of $\mathbf{D}$. Let $(x,n,nt),(y,s,sl)$ be two points in $I_{p}(M_{Y})$. Assume that $n\geq s$. 
\begin{align*}
d(\mathbf{D}(x,n,nt),\mathbf{D}(y,s,sl)) \leq d(\mathbf{D}(x,n,nt),\mathbf{D}(\mathfrak{q}_{s}(x,n),sl)) + d(\mathbf{D}(\mathfrak{q}_{s}(x,n),sl),\mathbf{D}(y,s,sl))\\
=d(\mathfrak{q}_{(1-t)n + t\gamma(n)}(x,n), \mathfrak{q}_{(1-l)s + l\gamma(s)}\mathfrak{q}_{s}(x,n)) +d(\mathfrak{q}_{(1-l)s + l\gamma(s)}\mathfrak{q}_{s}(x,n), \mathfrak{q}_{(1-l)s + l\gamma(s)}(y,s))
\end{align*} 
The first distance is just the straight line of length
\begin{align*}
|(1-t)n +t \gamma(n)-(1-l)s - l\gamma(s)| &\leq 2|nt-sl| + 3|n-s|\\
&\leq 2|nt-sl| + 3d((x,n),(y,s))
\end{align*} 
The distance between the second two points is the distance between the projections of $x$ and $y$ to the height $(1-l)s + l\gamma(s)$, measured within the appropriate nerve $|\mathcal{U}|$. Since the bonding maps are $1$-Lipschitz, we have
\begin{align*}
d(\mathfrak{q}_{(1-l)s + l\gamma(s)}\mathfrak{q}_{s}(x,n), \mathfrak{q}_{(1-l)s + l\gamma(s)}(y,s)) \leq d(\mathfrak{q}_{s}(x,n),(y,s)) \leq d((x,n),(y,s))
\end{align*}
Therefore the homotopy $\mathbf{D}$ is $4$-Lipschitz. Consider the composition 
\begin{align*}
\xi'  \mathbf{D}\circ H: I_{p}(cX)&\rightarrow M_{X}\\
(nx,n,nt)&\mapsto \xi' \mathbf{D}(\psi i'(\gamma'(n) x, \gamma'(n)), p\psi i'(\gamma'(n) x, \gamma'(n)) t)
\end{align*}
where $H: I_{p}(cX)\rightarrow I_{p}(M_{Y})$ is the rescaled version of $\psi i' s_{\xi' \psi}\times \id$. $\xi'  \mathbf{D}\circ H$ is a homotopy from $\xi' \psi i' s_{\xi' \psi}$ to  $\xi' \mathbf{D}_{1} \psi i' s_{\xi' \psi}$. We now show that $H$ is Lipschitz. Recall that $(L_{\psi})_{h}$ is the Lipschitz constant of $\psi$ restricted to heights $[1,2^{h+1}]$ and that $\{m'_{h}\}_{h\in \mathbb{N}_0}$ is the sequence of constants associated to $\gamma'$.  We can assume that $\psi = \widetilde{is_{l}\mathfrak{a}R}$ is height-decreasing by choosing $l$ sufficiently large. Let $(nx,n,nt),(sy,s,sl)$ be two points in $c_{ 2^{m'_{h}}[2^h,2^{h+1}]}(X \times [0,1])$. We have that
\begin{align*}
& d(\psi i'(\gamma'(n) x,\gamma'(n)), \psi i'(\gamma'(s) y,\gamma'(s))) \leq a_{h}(L_{\psi})_{h} d((\gamma'(n) x,\gamma'(n)),  (\gamma'(s) y,\gamma'(s)))\\
&= \frac{a_{h}(L_{\psi})_{h}}{2^{m'_{h}}} d((nx,n), (sy,s))\\
& |p\psi i'(\gamma'(n) x, \gamma'(n)) t-p\psi i'(\gamma'(s) y, \gamma'(s)) l)|\\
&\leq |p\psi i'(\gamma'(n) x, \gamma'(n)) t-p\psi i'(\gamma'(n)x, \gamma'(n)) l|+|p\psi i'(\gamma'(n) x, \gamma'(n)) l-p\psi i'(\gamma'(s) y, \gamma'(s)) l|\\
&= |t-l||p\psi i'(\gamma'(n) x, \gamma'(n))| + |p\psi i'(\gamma'(n) x, \gamma'(n)) -p\psi i'(\gamma'(s) y, \gamma'(s))|\\
&\leq |t-l|\frac{n}{2^{m'_{h}}} + \frac{(L_{\psi})_{h}a_{h}}{2^{m'_{h}}} d((nx,n),(sy,s))\\
&\leq \frac{1}{2^{m'_{h}}} (|nt-sl| + |n-s|) + \frac{(L_{\psi})_{h}a_{h}}{2^{m'_{h}}} d((nx,n),(sy,s))
\end{align*}
A similar computation works for heights in the interval in $ [2^{h+m'_{h}+1}, 2^{h+m'_{h+1}+1}]$. This shows that $H$ restricted to $c_{ [2^{h+m'_{h}},2^{h+m'_{h+1}+1}] }(X \times [0,1])$ is $\frac{3a_{h}(L_{\psi})_{h}}{2^{m'_{h}}}$-Lipschitz.
The computation for heights in the interval $[2^{h-1+m'_{h-1}},2^{h+m'_{h}}]$ is analagous, replacing $h$ with $h-1$. \\

Since in the construction of $S_{\xi'\psi}$, we have maximised over the Lipschitz constants of each successive bounded set, and $\mathbf{D}$ is height-decreasing, the homotopy $\xi \mathbf{D} H$ is uniformly locally $12$-Lipschitz with respect to $\{c_{ [2^{h+m'_{h}},2^{h+m'_{h+1}+1}] }(X \times [0,1])\}_{h\in \mathbb{N}_{0}}$ . Hence it is globally coarse. Precomposing with the coarse map $R \times \id$, we obtain a coarse homotopy between $\xi' \psi S_{\xi'\psi}$ and $\xi' \mathbf{D}_{1} \psi S_{\xi' \psi}$. This was step 1.\\

For the second step, we construct a coarse homotopy between $\xi' \mathbf{D}_{1} \psi S_{\xi' \psi}$ and $\xi' S_{\xi'} \psi S_{\xi' \psi}$ as the composition $\xi' \circ \mathbf{G}\circ \mathbf{F} \circ (R\times \id)$. $\mathbf{G}$ is defined as
\begin{align*}
\mathbf{G}: M_{Y} \times [0,1] &\rightarrow M_{Y}\\
(x,n,t) &\mapsto (1-t) \mathbf{D}_{1}(x,n) + t i's_{\xi'} R (x,n) 
\end{align*}
Why does this make sense? Consider a point $(x,n)$ with $x\in int(\sigma)$ in some appropriate nerve $|\mathcal{U}|$. $\mathbf{D}_{1}(\sigma \times \{ n\}) = \mathfrak{q}_{\gamma(n)}(\sigma \times \{n\}) = \phi(\sigma) \times \{\gamma(n)\}$ is a simplex at height $\gamma(n)$ for the appropriate bonding map $\phi$.  The image of $i's_{\xi'} R(int(\sigma) \times \{n\})$ lies in the set 
\begin{align*}
\bigcap_{v \in \phi (\sigma)} \langle st(v) \rangle  \times \{\gamma(n)\}
\end{align*} 
Therefore, the straight line between $ \mathbf{D}_{1}(x,n) $ and $i's_{\xi'} R (x,n)$ makes sense. Note that this homotopy is very much non-continuous. We define $\mathbf{F}$ as the map
\begin{align*}
\mathbf{F}: I_{p}(cX) &\rightarrow M_{Y} \times [0,1]\\
(nx,n,nt) &\mapsto (\psi i'(\gamma'(n)x,\gamma'(n)), t)
\end{align*}
Let $(nx,n,nt),(sy,s,sl)$ be two points heights in the interval  $[2^{h+m'_{h}},2^{h+m'_{h+1}+1}]$. Recall that we have 
\begin{align*}
d(\psi i'(\gamma'(n) x,\gamma'(n)), \psi i'(\gamma'(s) y,\gamma'(s)))
\leq \frac{a_{h}(L_{\psi})_{h}}{2^{m'_{h}}} d((nx,n), (sy,s))\\ 
|t-l| = \frac{1}{n}|nt-nl| \leq \frac{1}{2^{m'_{h}+h}}(|nt-sl| + |n-s|)
\end{align*}

Now we have to show that for a set sufficiently small in $M_{Y}\times [0,1]$, the image under $\mathbf{G}$ is actually sufficiently small. Let $\varepsilon, \delta$ be very small constants, to be determined later. Suppose we have two points $(x,n,t),(y,s,l) \in M_{Y} \times [0,1]$ below height $2^{l'}$, contained within $U_{c_{l'}\varepsilon} \times I_{\varepsilon} \times I_{\delta}$, where $I_{\varepsilon}$ only intersects at most one a gluing slice on the boundary and $U_{c_{l'}\varepsilon}$ is a set of diameter $<c_{l'}\varepsilon$ (with respect to the standard metric) contained within a simplex. The reason why we can restrict to this case is that for an arbitrary set $V$ of diameter $<\varepsilon'$ (with respect to the path metric) in $M_{Y}$ where $c_{l'}\varepsilon'<\frac{1}{2}$,  any two points $(u,n),(w,s)$ in $V$ can be split into at most $5$ segments of the form we have assumed. \\

To see this, let $n\geq s$. First we assume that the shortest path between $(u,n),(w,s)$ consists of $\omega_{(u,n)}\ast \eta_{s}$ where $\omega_{(u,n)}$ is the straight line downards and $\eta_{s} \subset (U_{\varepsilon'} \times \{s\})$ is a geodesic at height $s$ between $\mathfrak{q}_{s}(u)$ and $w$. Split $\omega_{(u,n)}$ into two paths which intersect at most one gluing slice. Now for two points $x,y$ in $U_{\varepsilon'}$, we consider their image in $(|\mathcal{U}|, d_{std})$, which have distance $< c_{l'}\varepsilon'<\frac{1}{2}$. Let $x = \sum_{v \in vert |\mathcal{U}|} x_{v} [v]$ and $y =  \sum_{v \in vert |\mathcal{U}|} y_{v} [v]$. Define
\begin{align*}
Z = \{v\in vert|\mathcal{U}|\,|\, x_{v}\neq 0, y_{v} \neq 0\}
\end{align*}
 and $a = \sum_{v \in Z} x_{v}$ and $b = \sum_{v \in Z} y_{v}$. The distance between $x$ and $y$ is 
\begin{align*}
d_{std}(x,y) = (1-a) + (1-b) + \sum_{v \in Z} |x_{v}-y_{v}| < c_{l'}\varepsilon'
\end{align*} 
Define $z_{x} = \frac{1}{a}\sum_{v\in Z} x_{v} [v]$ and $z_{y} = \frac{1}{b} \sum_{v\in Z} y_{v} [v]$. We compute
\begin{align*}
d_{std}(x,z_{x}) &= (1-a) + \sum_{v\in Z} (\frac{1}{a}-1) x_{v} = 2(1-a) < 2c_{l'}\varepsilon'\\
d_{std}(y,z_{y}) &< 2c_{l'}\varepsilon'\\
a d_{std}(z_{x},z_{y}) &= \sum_{v \in Z} |x_{v}-\frac{a}{b} y_{v}|\leq \sum_{v \in Z} |x_{v}- y_{v}|+ |\frac{a}{b}-1| y_{v}\\
&= (\sum_{v \in Z} |x_{v}- y_{v}|)+ |a-b| \leq (\sum_{v \in Z} |x_{v}- y_{v}|)+ |1-a| + |1-b| < c_{l'}\varepsilon'\\
d_{std}(z_{x},z_{y}) &<\frac{c_{l'}\varepsilon'}{a} <\frac{c_{l'}\varepsilon'}{1-c_{l'}\varepsilon'} <2c_{l'}\varepsilon'
\end{align*}
Therefore $\eta_{s}$ can be split into three straight paths of length less than $2c_{l'}\varepsilon'$, each contained within a simplex. \\

The only other case is if $n,s\in (2^{h},2^{h}+1]$ and the shortest path between $(u,n),(w,s)$ is a concatenation $\omega_{(u,n)} \ast \eta_{2^h} \ast \omega_{(w,s)}^{-1}$, where $\eta_{2^h}$ is a geodesic at height $2^h$. As before, we can split $\eta_{2^h}$ into three paths of the desired form.\\

Let us return to the proof. Recall we have two points $(x,n,t),(y,s,l) \in M_{Y} \times [0,1]$ below height $2^{l'}$, contained within $U_{c_{l'}\varepsilon} \times I_{\varepsilon} \times I_{\delta}$. Assume that $n\geq s$. We consider the points $(x,n,t),(x,s,t),(x,s,l),(y,s,l)$ under the image of $\mathbf{G}$ with respect to the standard metric $d_{std}$:
\begin{enumerate}
\item $d_{std}(\mathbf{G}(x,s,t), \mathbf{G}(x,s,l))$: Since the homotopy is the straight line between two points that lie within a simplex, it is $1$-Lipschitz in the time variable. Therefore $d_{std}(\mathbf{G}(x,s,t), \mathbf{G}(x,s,l))\leq |t-l|<\delta$. 
\item $d_{std}(\mathbf{G}(x,s,l), \mathbf{G}(y,s,l))$: We know that $d_{std}(\mathbf{D}_{1}(x,s), \mathbf{D}_{1}(y,s))\leq d_{std}(x,y)< c_{l}'\varepsilon$. Let $\sigma$ be the simplex that contains $x,y$. We know from previous calculations that $R(\sigma)$ has diameter less than $\frac{69}{4} = C$ (a conservative upper bound). Recall that $i's_{\xi'}$ is $\frac{1}{2^{b_{h'}}}$-Lipschitz for some $h'<l'$ where $2^{b_{h'}}$ is sufficiently large so that $\xi' i's_{\xi'}$ is uniformly locally $1$-Lipschitz. We can assume that $2^{b_{h'}}$ is also greater than the local Lipschitz constant $(L_{\xi'})_{h'}$ multipled by the Lipschitz constant $c_{h'}d_{h'}$ which we obtain from changing between metrics. Therefore $d(i's_{\xi'}R(x,s), i's_{\xi'}R(y,s))\leq \frac{C}{2^{b_{h'}}}$. In standard metric, $d_{std}(\mathbf{G}(x,s,l), \mathbf{G}(y,s,l))\leq \max\{\frac{Cc_{h'}}{2^{b_{h'}}},c_{l}'\varepsilon\}$. To see this, observe that for four points $v_1,v_2,w_1,w_2$ in Euclidean space, 
\begin{align*}
\|(1-t)v_1+tw_1 - (1-t)v_2 + tw_2\| \\
\leq (1-t)\|v_1-v_2\| + t\|w_1-w_2\| \leq \max \{\|v_1-v_2\|, \|w_1-w_2\|\}
\end{align*}
\item $d_{std}(\mathbf{G}(x,n,t), \mathbf{G}(x,s,t))$:  Ignore first the possible discontinuity related to $R$ at the boundary. We have that $R(x,n)$ and $R(x,s)$ get sent to the same $x$-coordinate, so their distance in $cY$ is less than $2|n-s|$. Remember that by construction, $i' s_{\xi'}$ is actually $\frac{1}{2^{b_{h'}}}$-Lipschitz with respect to the product metric on each $\tilde{M}^{h'}$ before gluing.  Because the shrinking maps are constructed using powers of $2$, the image of $s_{\xi'}(x,n)$ and $s_{\xi'} (x,s)$ do not intersect a gluing component except possibly at the boundary - we are allowed to compute distances in one such $\tilde{M}^{h'}$. Therefore, the $|\mathcal{U}_{h'}|$-component of $i's_{\xi'}R(x,n)$ and $i's_{\xi'} R(x,s)$ have distance less than than $2\cdot \frac{1}{2^{b_{h'}}} |n-s|$. Since $\mathbf{D}_1(x,n)$ and $\mathbf{D}_{1}(x,s)$ have the same $x$-coordinate, by the previous case we conclude that 
\begin{align*}
d_{std}(\mathbf{G}(x,n,t), \mathbf{G}(x,s,t))\leq (2\cdot \frac{c_{h'}}{2^{b_{h'}}}+1) |n-s|
\end{align*}
If we do not intersect a gluing slice, this finishes the estimate. If $n$ or $s$ lie on a gluing slice there is a discontinuity to consider. By the previous case, the discontinuity is bounded by $\frac{Cc_{h'}}{2^{b_{h'}}}$, so in total we get
\begin{align*}
d_{std}(\mathbf{G}(x,n,t), \mathbf{G}(x,s,t))< \frac{Cc_{h'}}{2^{b_{h'}}} + (2\cdot \frac{c_{h'}}{2^{b_{h'}}}+1) \varepsilon <\frac{Cc_{h'}}{2^{b_{h'}}} + 3\varepsilon
\end{align*}  
in both cases.
\end{enumerate} 

Therefore, $d_{std}(\mathbf{G}(x,n,t),\mathbf{G}(y,s,l))<  \delta +\max\{\frac{Cc_{h'}}{2^{b_{h'}}},c_{l}'\varepsilon\}+\frac{Cc_{h'}}{2^{b_{h'}}} + 3\varepsilon$. We now let $\varepsilon =  \frac{a_{h}(L_{\psi})_{h}}{2^{m'_{h}}}$ and $\delta = \frac{1}{2^{m'_{h}+h}}<\varepsilon$, which gives us $d_{std}(\mathbf{G}(x,n,t),\mathbf{G}(y,s,l))<  (5c_{l'}\varepsilon + 2\frac{Cc_{h'}}{2^{b_{h'}}})$. Recall that $2^{m_{h}'}> a_{h} (L_{\psi})_{h} (L_{\xi'})_{l'}c_{l'}d_{l'}$ and that $2^{b_{h'}} > c_{h'}d_{h'} (L_{\xi'})_{h'}$.  Postcomposing with the locally Lipschitz map $\xi'$, we obtain 
\begin{align*}
d_{path}(\xi'\mathbf{G}(x,n,t),\xi'\mathbf{G}(y,s,l))&\leq (L_{\xi'})_{h'} d_{h'} (5c_{l'}\varepsilon + 2\frac{Cc_{h'}}{2^{b_{h'}}})\\
&\leq (L_{\xi'})_{l'}d_{l'}\frac{5c_{l'}a_{h}(L_{\psi})_{h}}{2^{m'_{h}}} +2(L_{\xi'})_{h'} d_{h'}\frac{Cc_{h'}}{2^{b_{h'}}}\\
&<5 +2C
\end{align*}
Now we finally show that $\xi\circ \mathbf{G}\circ \mathbf{F}$ is a coarse homotopy. Consider $(nx,n,nt),(sy,s,sl)$, two points in $I_{p}(cX)$ with heights in the interval $ [2^{h+m'_{h}},2^{h+m'_{h+1}+1}]$ and distance less than $Q\in \mathbb{N}$. $\mathbf{F}(nx,n,nt), \mathbf{F}(sy,s,sl)$ live below height $2^{l'}$ and are contained in $V_{Q\varepsilon} \times I_{Q\delta}$. Therefore, the shortest path between them can be split up into $2Q$ segments each lying within $V_{\frac{\varepsilon}{2}} \times I_{\delta}$. By previous discussion we can split this further to obtain $10Q$ path segments each lying within a set of the form $U_{c_{l'}\varepsilon} \times I_{\varepsilon} \times I_{\delta}$ with respect to the standard metric. The endpoints of each of these segments, after composing with $\xi' \mathbf{G}$, are distance $< 5+2C$ apart. We obtain
\begin{align*}
d_{path}(\xi'\mathbf{GF}(nx,n,nt),\xi'\mathbf{GF}(sy,s,sl))&\leq 10Q(5+2C)
\end{align*}
Therefore, a set of diameter $Q\in \mathbb{N}$ in $c_{[2^{h+m'_{h}},2^{h+m'_{h+1}+1}]}(X \times [0,1])$ gets sent to a set of diameter $(50+20C)Q$ under the map $\xi' \mathbf{GF}$. By a similar argument to Lemma \ref{shrinkcoarse}, for arbitrary points $(nx,n,nt),(sy,s,sl)$ in $I_{p}(cX)$ and $Q>0$, there exists a constant $T_{Q}$ such that 
\begin{align*}
d((nx,n,nt),(sy,s,sl))<Q \\
\implies d(\xi' \mathbf{GF}(nx,n,nt),\xi'  \mathbf{GF}(sy,s,sl)) < \max\{3(50+20C)Q, (50+20C)Q+T_{Q}\}
\end{align*} This shows that $\xi' \mathbf{GF}$ is coarse. By precomposing with the coarse map $R \times \id$, we obtain that $\xi' \mathbf{D}_{1} \psi S_{\xi' \psi}$ and $\xi' S_{\xi'} \psi S_{\xi' \psi}$ are coarsely homotopic. \\

The two steps show that $\xi' \psi S_{\xi'\psi}\simeq \xi' \mathbf{D}_{1} \psi S_{\xi' \psi}\simeq \xi' S_{\xi'} \psi S_{\xi' \psi}$ which was the required homotopy to conclude 
\begin{align*}
\xi' S_{\xi'} (is_{l}\mathfrak{a}R) \simeq \id_{M_{X}}
\end{align*}

Therefore, $\xi' S_{\xi'}$ is a coarse homotopy inverse to $is_{l}\mathfrak{a}R$. 
\end{proof}

\end{document}
